\phantomsection\bheading{section}{3.\quad The inference algorithm}{section}{\protect\numberline{3.}The inference algorithm}

\phantomsection\bheading{subsection}{3.1\quad Shape}{subsection}{\protect\numberline{3.1}Shape}

Four passes per module, after type checking, in the slot \code{Effects.run} occupies today (\code{checker-v2.md} §26: after P5, before P6, with every solved type available; the backend “sees no types”, §13.4, so nothing later could do this):

\begin{enumerate}
\item \textbf{Local flow} (§3.2): one forward walk of each declaration's body in evaluation order, computing $A(e)$ for every list-typed expression. Calls use callees' summaries; a callee in the same strongly connected component uses the current round's summary.
\item \textbf{Liveness} (§3.3): one backward walk per declaration, computing $\mathit{live}(h,\allowbreak{}p)$ for the holders the forward walk named.
\item \textbf{Check and summarise} (§3.4): the K-rules at every demand; the declaration's summary read off the walks.
\item \textbf{The group fixpoint} (§3.5): repeat 1–3 for a recursive group until no summary changes, then validate.
\end{enumerate}

Function-typed values carry \textbf{ownership classes} (§3.6), solved with $\sqsubseteq$ edges in the same pass. Imported summaries come from interfaces (§3.9). The pass runs on every build, like the effect pass; it produces diagnostics, unlike the write-set pass.

\phantomsection\bheading{subsection}{3.2\quad Local flow: directional, flow-sensitive, in evaluation order}{subsection}{\protect\numberline{3.2}Local flow: directional, flow-sensitive, in evaluation order}

BIR is already a sequence of instructions in evaluation order with \code{let} binding and \code{case} branching (\code{frontend.md} §3.6; \code{language.md} §6, \emph{Evaluation order}). The walk keeps $\Sigma$ and applies §2.5's rules; a \code{case} forks $\Sigma$ per arm and joins the results by union. There is no loop inside a body except the self tail call, which the walk handles as a recursive call of the same declaration (so the parameters' cells on re-entry are the union of the entry cells and the jump's arguments — one fixpoint over the body, which converges in at most as many rounds as there are distinct cells, in practice one or two).

Why directional and flow-sensitive both matter, in one example:

\begin{Verbatim}
xs = [ 1, 2, 3 ]            -- A(xs) = {site₁}
ys = xs                     -- A(ys) = {site₁}: a copy of the set, not a class
zs = List.set ys 0 9        -- demand: site₁ unique?  holders: xs, ys. xs dead after? …
List.length zs              -- …yes: xs is never read again. Accepted.
\end{Verbatim}

A unification-based model would put \code{xs}, \code{ys}, \code{zs} in one class and then any later use of any of them would taint all three; here \code{zs} is a new holder of the \emph{same} cell (ownership transferred), and \code{xs} is dead.

The cost of pass 1 is linear in the instructions times the width of the abstract values (the number of list-typed k-paths under a type, bounded by the type; the cell sets are usually singletons). Abstract values are hash-consed as \code{write-sets.md} §3.1 does, so a chain of \code{if}s does not duplicate them.

\phantomsection\bheading{subsection}{3.3\quad Liveness}{subsection}{\protect\numberline{3.3}Liveness}

A standard backward pass over the same instruction sequence, with a bit per (variable, path) and the uses of §2.7. Branches join by union (a holder live in either arm is live before the \code{case}). Closures: a captured path is live \textbf{wherever the closure value is live} — at every call, store, return or pass of the closure, which the walk sees as ordinary uses of the closure's binding — and \textbf{from the creation onward, for ever}, when the closure escapes (passed to an $\mathit{escapes}=\mathsf{yes}$ position, stored, returned). So \code{g = λ\_\allowbreak{} → List.\allowbreak{}length xs; ys = List.\allowbreak{}push xs 1; \{ m | f = g \}} is rejected: $g$ is read after the push and holds \code{xs}. The cost is linear.

\phantomsection\bheading{subsection}{3.4\quad The check, and reading off the summary}{subsection}{\protect\numberline{3.4}The check, and reading off the summary}

At each demand apply K1–K4. Then:

\begin{itemize}
\item $\mathit{use}(i,\allowbreak{}q)=\mathsf{consumed}$ if some demand's operand may hold $\pi _{i}.q$ (K3 was discharged for every \emph{local} holder; the caller discharges the rest);
\item $\mathit{use}(i,\allowbreak{}q)=\mathsf{shared}$ if $\pi _{i}.q \in A(\mathit{result})@q'$ for some $q'$, or $\pi _{i}.q$ was stored in an escaping place, or passed to a $\mathsf{shared}$ position of a callee, or captured by an escaping closure;
\item $\mathit{use}(i,\allowbreak{}q)=\mathsf{borrowed}$ otherwise;
\item $\mathit{res}(q')=\{\pi _{i}.q\mid \pi _{i}.q\in A(\mathit{result})@q'\}\cup \{\mathsf{fresh}\mid \mathit{some}\ \mathit{site}\in A(\mathit{result})@q'\}$;
\item $\mathit{fun}(j)$ from the calls of $\pi _{j}$ the walk met: $\mathit{calls}$ by counting on paths (any call inside a tail-call loop or passed to a $\mathsf{many}$ position is $\mathsf{many}$), $\mathit{escapes}$ from the holders, $\mathit{supply}$ from the K-check at each call (owned iff the argument's cell was unique there), $\mathit{need}$ from whether the result's cells reach $A(\mathit{result})$ or a store.
\end{itemize}

$\mathsf{consumed}$ and $\mathsf{shared}$ are not ordered; a parameter can be both ($\mathsf{consumed}$ and in $\mathit{res}$, which is the normal writer shape). The three-valued $\mathit{use}$ is the aspect lattice of Aspinall, Hofmann and Konečný — “Variables may only be accessed once with aspect 1. Variables can be used many times with aspect 2 … Finally, variables can be freely used with aspect 3” (2008 §1, p. 4) — with “once” replaced by “unique at the site”.

\phantomsection\bheading{subsection}{3.5\quad Recursion: the optimistic fixpoint, and why it is the best answer}{subsection}{\protect\numberline{3.5}Recursion: the optimistic fixpoint, and why it is the best answer}

A strongly connected component of the call graph (within one module; the module graph is a DAG, so cross-module recursion cannot occur) is solved together. Every member starts at the most optimistic summary: every parameter $\mathsf{borrowed}$, every result $\{\mathsf{fresh}\}$, every function-typed parameter $\mathit{calls}\le 1,\allowbreak{}\mathit{escapes}\ \mathsf{no}$. The body is analysed with those; the summaries that come out are joined with the assumption (pointwise: $\mathsf{borrowed}\sqsubset \mathsf{shared}$, $\mathsf{borrowed}\sqsubset \mathsf{consumed}$, $\mathsf{fresh}$ $\cup$ parameters; $\le 1\sqsubset \mathsf{many}$, $\mathsf{no}\sqsubset \mathsf{yes}$); repeat until nothing changes. The lattice is finite (two independent bits per parameter path, a set per result path, two bits per function parameter) and every transfer is monotone (more sharing in a callee's summary makes more holders and more escapes in the caller, never fewer), so the iteration terminates in at most $2\cdot\lvert\mathit{paths}\rvert + \lvert\mathit{results}\rvert + 2\cdot\lvert\mathit{fun\ params}\rvert$ rounds and in practice one or two. \textbf{Then the group is validated}: each body is checked once more against the final summaries, and every K-rule must pass; a summary that reached a fixpoint is a consistent \emph{typing} of the group, and §4.3's Lemma 3 argues its soundness by induction on the number of calls an execution makes.

This is the procedure the literature agrees on. Aspinall, Hofmann and Konečný: “for every program typable in our system, there is a typing with best (i.e. largest) usage aspects. These usage aspects can be automatically reconstructed using an iterative search for a fixed point, starting with the most optimistic typing of all functions in the program” (2008 §6.2, p. 36 — the algorithm itself is in Konečný 2003, not in the library). Lean: “we infer $\beta(c)$ by starting with the approximation $\beta(c) = B^n$, then we compute $S = \mathit{collect}_O(b)$, update $\beta(c)_i := O$ if $y_i \in S$, and repeat the process until we reach a fix point” (Ullrich \& de Moura 2019 §5.2, p. 6). Brandon et al.: “each group of equations in Figure 8 defines a monotone function … Existence of a least fixed point for each group of equations is justified by Knaster–Tarski, and termination is justified by finiteness. In other words, our equations can be viewed as a Datalog program” (2026 §4.4, p. 16). The write-set pass does the same for its summaries (\code{write-sets.md} §4.3, Kleene iteration from $\bot$). Clarke et al.'s complaint that ownership inference “cannot give a best solution” (research 69 §1.1) does not apply: here the constraints are inclusions with a least solution, and the optimistic start finds it.

Where Clean gave up — “we do not try any specific instances but consider the expression untypable” (Barendsen \& Smetsers 1996 §8, p. 33) and “the type of every binding in a recursive binding group must be non-unique” (de Vries 2007 §7, p. 12) — this checker does not, because its unknowns are facts about flows, not attributes on types that a rank restriction on arrows prevents from being generalised.

\phantomsection\bheading{subsection}{3.6\quad Function values: ownership classes on function types}{subsection}{\protect\numberline{3.6}Function values: ownership classes on function types}

A function-typed parameter $g$ of $f$ has no body to analyse. Two ways exist: assume the worst (every argument to $g$ shared, $g$'s result $\top$, $\mathit{calls}\ \mathsf{many}$) — Huisinga's “higher-order functions are always shared” (2023 §4.1, p. 47) and the Lean thesis's “considerable limitation” — or make $f$'s summary \emph{polymorphic} in $g$'s. beni already does the second for effect bits: “a reference to a top-level declaration instantiates its summary … each class of the scheme gets a fresh class at the use, which gets the scheme class's rung and $d\sqsubseteq c$ for every scheme class $d$ that $c$ depends on” (\code{transparent-effects-proposal.\allowbreak{}md} §14.3 rule 3), which is “what makes \code{List.map} serve a pure and a suspending callback with one definition”. Ownership takes the same seat:

\begin{itemize}
\item every function type in a scheme carries an \textbf{ownership class}: $\mathit{fun}(j)$'s four facts and, for its own parameters and result, $\mathit{use}$/$\mathit{res}$ variables;
\item $f$'s summary is expressed over its function-typed parameters' classes: \code{map}'s result is $\mathsf{fresh}$; its callback is $\mathit{calls}\ \mathsf{many},\allowbreak{}\mathit{escapes}\ \mathsf{no},\allowbreak{}\mathit{supply}\ \mathsf{borrowed}(\mathit{the}\ \mathit{element}),\allowbreak{}\mathit{need}\ \mathsf{fresh}$; \code{foldl : List a [borrowed], b [consumed], (a, b [consumed] → b) [calls many, supply: borrowed, owned; need: owned] → b [= π₂ ∪ g.\allowbreak{}res]};
\item at a call \code{f … λ …} the lambda's actual summary is related to the instantiated class by \textbf{directional edges}, not unification: the lambda may consume a parameter only where $f$ supplies an owned cell ($\mathit{need}_{\lambda}(k)=\mathsf{consumed}\Longrightarrow \mathit{supply}_{f}(k)=\mathsf{owned}$); the lambda's $\mathit{once}$ requires $\mathit{calls}\le 1$; the lambda's result aliases flow into $f$'s $\mathit{need}$; the lambda's $\mathsf{shared}$ captures flow into $f$'s result if $f$ returns $g$'s result; $\mathit{escapes}=\mathsf{yes}$ escapes every capture. A violated edge is an error at the call, naming the lambda and the parameter. This is subtyping, “not just unification”, exactly §4.4 of the effects proposal.
\end{itemize}

The classes ride on the solved types, so a function value that reaches $f$ through a record field or a parameter of a parameter is handled by the same instantiation (the class is where the function type is). Rank-1 is enough: Lorenzen et al. note that “higher-order functions in Rust that pass newly borrowed values to their callbacks must have higher-rank types … while in our system such functions have rank-1 types that can be fully inferred” (2024 §8.3, p. 26), and beni's arrows are n-ary and saturated, so there is no spine to propagate along.

What this does \textbf{not} do is infer, inside $f$, that a particular closure \emph{only reads} a capture and may therefore be called many times while the capture is later consumed. That fact is the lambda's ($\mathit{cap}\ z=\mathsf{borrowed}$), decided where the lambda is written, and it composes: \code{List.\allowbreak{}map xs (λx → f x ys)} has $\mathit{cap}\ \mathit{ys}=\mathsf{borrowed}$, \code{map}'s $\mathit{calls}=\mathsf{many}$ is fine for a closure that consumes nothing, the closure does not escape, so \code{ys} is live only during the \code{map} call, and a later \code{List.set ys …} is accepted. Deng et al. 2025 get the same by effects ($\{u:c\}$ for a reader, $\{u:c;k:c\}$ for a consumer, Fig. 6, p. 11) and Capybara by capture modes (a read-only capturing closure “never blocks a later write”, ms.tex:647–671); both declare, this infers, because the lambda is in front of the checker.

\phantomsection\bheading{subsection}{3.7\quad \code{where} clauses and generic code}{subsection}{\protect\numberline{3.7}\code{where} clauses and generic code}

A method call \code{x.m a} inside a function with a \code{where} clause calls evidence the caller supplies. Two cases:

\begin{itemize}
\item the well-known \code{eq} and \code{compare} are \code{sync} and read-only by contract (\code{boundary.md} §4): summary $\mathsf{borrowed},\allowbreak{}\mathsf{borrowed}\to \mathit{crosses}$. No loss.
\item a user method \code{a.\allowbreak{}update : a, Int → a}: its summary is a \textbf{class} on the constraint's method type, instantiated at each call site from the dispatch table's resolved term (\code{checker-v2.md} §13.1: \code{top}, \code{ext}, \code{derived}, \code{param}), exactly as its effect class is. Inside the generic body the class's variables are what the summary is written over; at the use they are bound.
\end{itemize}

So generic code loses nothing by itself. What it cannot do is write a type variable's value in place: \code{f : a → a} has no demand to call. A generic container function — \code{Dict.\allowbreak{}update d k (λv → List.\allowbreak{}push v x)} — is the container's business (§5.4), and its callback's $\mathit{need}$ is a class the caller's lambda is checked against.

Parametricity supplies the result aliasing for free: a function of type \code{a → a} can only return its argument (or something a callback gave it), so \code{identity}'s summary is $\mathit{res}=\{\pi _{1}\}$ without analysis — Futhark's observation that “the only way this function can return an $a$ is by returning the $a$ we give it” (2026 aliasing post, ll. 555–612) and the $\mathit{id}(x:T^{\blacklozenge}):T^{x}$ of Wei et al. The checker gets it from the body anyway (\code{identity x = x} analyses to that), and from the type for a \code{foreign} over a type variable (\code{Debug.log}'s row in \code{write-sets.md} §3.7) \textbf{only when no \code{foreign type} with a tracked parameter stands between the argument and the result}: \code{Task.\allowbreak{}join : Fiber a → a} returns a value the runtime retains and hands to every joiner, so its result is $\top$ (§2.2's runtime-retained row), not “fresh from nowhere”.

\phantomsection\bheading{subsection}{3.8\quad Complexity, and the budget}{subsection}{\protect\numberline{3.8}Complexity, and the budget}

Per declaration: pass 1 is $O(\lvert\mathit{body}\rvert\cdot w)$ where $w$ is the number of tracked paths of the widest type the body mentions (bounded by k and the type; for a model record with three list fields $w=3$); pass 2 is $O(\lvert\mathit{body}\rvert\cdot w)$; the check is $O(\mathit{demands}\cdot\mathit{holders})$, which is quadratic \emph{within one declaration} when a cell has many alias bindings — \code{update} is the declaration that has both, and §9 measures it. A recursive group of $n$ declarations (top-level, or the \code{let} functions of one body, which form their own component) runs $\le$ $h$ rounds where $h \le 2\cdot P + R + 2\cdot F$ is the lattice height over the group's parameter paths, result paths and function parameters, so the group costs $O(h\cdot\Sigma\lvert\mathit{body}\rvert\cdot w)$. Interfaces add a few words per scheme site. Nothing is quadratic in the module, and nothing is whole-program: summaries are read from interfaces exactly as effect summaries are (\code{checker-v2.md} §26, “imported summaries are applied at instantiation”).

\code{fast-compiler.md} §2 asks for > 250k LOC/s per core for checking; the effect inference landed inside the plan's “$\le$ +10 \% of check” bound (§14.4 there), and this pass has the same shape with a wider per-site record. The honest unknown is $w$ on wide model records and the holder count at a demand inside a large \code{update}; §9 measures both. A cap, if one is needed, goes on $w$ (a type with more than, say, 64 tracked paths is tracked as a whole) and yields $\top$ — a \emph{limit} report, never a wrong answer, the shape every cap in \code{write-sets.md} §6.1 has.

\phantomsection\bheading{subsection}{3.9\quad Summaries in the interface, and the firewall}{subsection}{\protect\numberline{3.9}Summaries in the interface, and the firewall}

The summary of every \code{pub} declaration is published in the interface record as a block beside the effect block (\code{transparent-effects-proposal.\allowbreak{}md} §14.6): per scheme site (found by the same root/path encoding), the $\mathit{use}$ and $\mathit{res}$ facts, and per function-type site its class's four facts. The block is \textbf{hashed with the record} (\code{fast-compiler.md} §8.1), so a body edit that changes a summary changes the hash and the importers are re-checked, which is the only correct behaviour: Klabnik's “if it did [infer signatures], changing the body of the function would change its signature” (research 69 §1.3 item 5) is a description of this design, and the cost is churn, which research 19 §4 priced for inferred \code{where} suffixes at 4.4–6.5× the churn of annotated code. §9 measures it for summaries; §10 offers an optional annotation that freezes one.

The cache key does not change: the block is a function of the module's source and its imports' interfaces, which the key already pins (\code{checker-v2.md} §28, same argument for the boundary rows).

\phantomsection\bheading{subsection}{3.10\quad Determinism}{subsection}{\protect\numberline{3.10}Determinism}

Every order the pass uses is the checker's: SCC order tie-broken by source, instruction order within a declaration, name text for fields and paths, store-variable index for classes (\code{checker-v2.md} §17). Cells are named by instruction index and parameter position. No hash of a pointer, no thread order. The \code{--jobs=1} versus \code{--jobs=8} determinism scenario gains the ownership dump (§9.1) beside the dispatch and writes dumps.

\phantomsection\bheading{subsection}{3.11\quad What the pass is not}{subsection}{\protect\numberline{3.11}What the pass is not}

It is not an optimisation that decides per site whether to write in place: under the rule every accepted demand is in place, so the backend needs no per-site bit (§7.5). It is not the optimiser's analysis: its verdict is a function of the source and the rule set, and inlining, specialisation or dead-binding elimination never change it — the optimiser is instead \emph{told} where the demands are and moves no read across one (A5, §7.6) — McCall's condition that the guaranteed set be “real guarantees … even in debug builds” and never be “defined … by what \code{-O} is able to eliminate in a particular release” (Swift forums 2025, \#16).
